% FORMALIZACION_NUEVA de la realización homogénea preexistente.
% Prioridad de la fórmula integrada: procedencia histórica pendiente de cotejo.
% Conserva la ecuación y el dominio del capítulo de rebote; no selecciona s_0.
\section{Solución homogénea explícita y completitud causal}
\label{vii:sec:solucion-homogenea}
La determinación de la amplitud permite desarrollar una realización
completa del rebote, además del argumento local. Se conserva la coordenatización
espacialmente plana de \eqref{vii:eq:dinamica-homogenea}, en unidades
\(c=1\), con \(\rho_0,s_0,G>0\), y se especializa el sector material
a polvo, \(w=0\), con \(\Lambda_{\mathrm{eff}}=0\).
El coeficiente \(s_0\) es el obtenido del estado
material en la realización correspondiente; la integración siguiente
emplea ese coeficiente y conserva su dominio.

\begin{theorem}[Integración exacta del rebote de polvo]
\label{vii:thm:rebote-polvo}
Pongamos
\[
 a_{\rm b}^3=\frac{s_0}{\rho_0},\qquad
 \tau^2=\frac{s_0}{6\pi G\rho_0^2},\qquad u=t-t_{\rm b}.
\]
La solución del sistema homogéneo con
\(a(t_{\rm b})=a_{\rm b}\) y \(H(t_{\rm b})=0\) es, para todo
\(t\in\mathbb R\),
\begin{equation}
 a(t)=a_{\rm b}\left(1+\frac{u^2}{\tau^2}\right)^{1/3},
 \qquad H(t)=\frac{2u}{3(\tau^2+u^2)},\qquad
 \dot H(t)=\frac{2(\tau^2-u^2)}{3(\tau^2+u^2)^2}.
 \label{vii:eq:solucion-polvo}
\end{equation}
Tiene un único mínimo, en \(t_{\rm b}\); es contractiva antes y
expansiva después. La evolución es suave y satisface
\(a(t)\geq a_{\rm b}>0\).
\end{theorem}
\begin{proof}
Sea \(y=a^3\). La ecuación de Friedmann da
\[
 \dot y^{\,2}=24\pi G(\rho_0y-s_0).
\]
La función \(y=s_0/\rho_0+6\pi G\rho_0u^2\) verifica esa
identidad y las condiciones iniciales. Sus derivadas dan las
tres expresiones de \eqref{vii:eq:solucion-polvo}. Para comprobar
también la ecuación de evolución, escribamos
\[
 \rho_{\rm b}=\frac{\rho_0^2}{s_0},\qquad
 x=\frac{u}{\tau}.
\]
Entonces
\[
 \rho_{\rm m}=\frac{\rho_{\rm b}}{1+x^2},\qquad
 \rho_{\rm s}=\frac{\rho_{\rm b}}{(1+x^2)^2},\qquad
 4\pi G\rho_{\rm b}=\frac{2}{3\tau^2}.
\]
Sustituir en \(\dot H=-4\pi G(\rho_{\rm m}-2\rho_{\rm s})\)
da exactamente la tercera fórmula. Se satisface, por tanto,
el sistema completo \((\dot a,\dot H)\), cuyo campo es suave
para \(a>0\). Su unicidad local, prolongada sobre la solución
explícita, fija la evolución por esas condiciones iniciales.
El denominador es positivo para todo \(t\), y las afirmaciones
sobre el mínimo y el signo siguen de \(H\).
\end{proof}

\begin{figure}[htbp]
\centering
\begin{tikzpicture}[x=1.45cm,y=1.15cm,>=Latex]
\draw[->] (-3.25,0)--(3.3,0) node[right] {\(\displaystyle u/\tau\)};
\draw[->] (0,0)--(0,2.8) node[above] {\(\displaystyle a/a_{\rm b}\)};
\draw[densely dashed,gray] (-3.1,1)--(3.1,1);
\draw[very thick,ink,domain=-3:3,samples=121,smooth]
 plot (\x,{(1+\x*\x)^(1/3)});
\fill[accent] (0,1) circle (1.6pt);
\node[anchor=north east] at (-.08,.98) {\(1\)};
\node[ink] at (-2,2.65) {Contracción};
\node[ink] at (2,2.65) {Expansión};
\foreach \t in {-3,-2,-1,1,2,3}
{\draw (\t,.045)--(\t,-.045) node[below] {\(\t\)};}
\end{tikzpicture}
\caption{Perfil adimensional de la solución demostrada:
\(a/a_{\rm b}=(1+(u/\tau)^2)^{1/3}\).
Los parámetros \(a_{\rm b}\) y \(\tau\) se calculan desde
\(\rho_0,s_0,G\); el dibujo representa toda la colección reescalada.}
\label{vii:fig:rebote-polvo}
\end{figure}

\subsection{Curvatura finita y extensión de las geodésicas}
La métrica realizada en esta coordenatización es
\(g=-dt^2+a(t)^2\,d\mathbf x^2\) sobre
\(\mathbb R\times\mathbb R^3\). Con la convención en que
el escalar de curvatura plano es \(R=6(\dot H+2H^2)\), la solución da
\[
 R(t)=\frac{4\tau^2+\frac43u^2}{(\tau^2+u^2)^2}.
\]
En particular, \(R(t_{\rm b})=4/\tau^2\) es finito.
En un marco ortonormal, las componentes de Riemann dependen de
\(H^2\) y de \(\dot H+H^2\), que son funciones acotadas en
\(\mathbb R\). Lo mismo ocurre con las contracciones polinómicas
de esas componentes; por ejemplo,
\[
 R_{\mu\nu\rho\sigma}R^{\mu\nu\rho\sigma}
 =12\bigl[(\dot H+H^2)^2+H^4\bigr].
\]

\begin{proposition}[Completitud de las geodésicas causales]
\label{vii:prop:completitud-causal-polvo}
La realización métrica precedente es geodésicamente completa
para las geodésicas temporales y nulas de su conexión de
Levi--Civita, en ambas direcciones.
\end{proposition}
\begin{proof}
Las traslaciones espaciales proporcionan las constantes
\(p_i=a^2\,dx^i/d\lambda\). La normalización causal da
\[
 \left|\frac{dt}{d\lambda}\right|
 =\sqrt{\epsilon+\frac{|\mathbf p|^2}{a(t)^2}},
 \qquad \epsilon=1\ \text{en el caso temporal},\quad
 \epsilon=0\ \text{en el caso nulo}.
\]
Una geodésica nula no constante tiene \(|\mathbf p|>0\).
Como \(a\geq a_{\rm b}\), en el caso temporal
\[
 \left|\frac{d\lambda}{dt}\right|
 \geq\left(1+\frac{|\mathbf p|^2}{a_{\rm b}^2}\right)^{-1/2}>0,
\]
y en el nulo
\(\left|d\lambda/dt\right|=a/|\mathbf p|
 \geq a_{\rm b}/|\mathbf p|>0\).
Así, alcanzar \(t=\pm\infty\) requiere longitud afín infinita.
En un intervalo temporal compacto, \(a\) y sus derivadas son
suaves y \(a\) está separado de cero. Las ecuaciones explícitas
anteriores y \(dx^i/d\lambda=p_i/a^2\) permiten continuar la
geodésica a través de sus extremos finitos. Quedan cubiertas
las dos posibilidades de prolongación.
\end{proof}

El resultado de completitud pertenece a esta realización homogénea de
polvo, con sus condiciones de espacio, materia y amplitud. Su
persistencia local bajo perturbaciones del umbral conserva las cotas
del teorema~\ref{vii:thm:persistencia}; esa propiedad local y la
completitud global de la solución explícita son conclusiones distintas.
La escala mínima procede de la relación entre densidad material y
respuesta torsional. La memoria interviene así en una evolución
geométrica verificable, con una realización regular de la contracción
y la expansión.
